\documentclass[12pt,a4paper]{article}

\usepackage[utf8]{inputenc}
\usepackage[T1]{fontenc}
\usepackage{lmodern}
\usepackage{charter}
\usepackage[ngerman]{babel}
\usepackage[margin=1.6cm]{geometry}

\usepackage{enumitem}
\usepackage{titlesec}
\usepackage[document]{ragged2e}
\usepackage[shortcuts]{extdash}

\hyphenpenalty=10000
\exhyphenpenalty=0
\emergencystretch=2em

\pagestyle{empty}

\setlist[itemize]{
  noitemsep,
  topsep=5pt,
  leftmargin=2em,
}

\setlist[description]{
  leftmargin=!,
  labelwidth=1.5cm,
  style=nextline
}

\titleformat{\section}{\fontsize{18}{22}\bfseries}{}{0em}{}[{\titlerule[1.5pt]}]
\titlespacing*{\section}{0pt}{18pt}{10pt}

\begin{document}

\noindent
{\fontsize{24}{28}\selectfont\bfseries Features} \\
\begin{minipage}[t]{0.7\textwidth}
  \vspace{8pt}
  {\itshape \textbf{Fach:} RESTFUL Web Services} \\
  {\itshape \textbf{Gruppe:} Blue Team}
\end{minipage}
\begin{minipage}[t]{0.3\textwidth}
  \flushright
  \vspace{8pt}
  \itshape \textbf{Datum:} \today
\end{minipage}

\section{User}
\subsection{Registrierung}
Benutzer können ein Konto registrieren, indem sie ihre grundlegenden 
Informationen angeben.

\subsection{Authentifizierung}

Bei jeder Anmeldung durchlaufen Benutzer einen Authentifizierungsschritt, 
um die Sicherheit ihres Kontos zu gewährleisten. Dieser zusätzliche Schritt 
dient dazu, unbefugten Zugriff zu verhindern und die Identität des Benutzers 
zu bestätigen.

\subsection{Profilverwaltung}
Benutzer können ein persönliches Profil erstellen und verwalten, das 
grundlegende Informationen wie Avatar, Benutzername, E-Mail-Adresse 
und Passwort enthält. Sie haben die Möglichkeit, ihre persönlichen Daten 
jederzeit einzusehen und zu bearbeiten.

\subsection{Kategorienverwaltung}

Benutzer können vordefinierte Kategorien verwenden, um schneller zu starten. 
Zusätzlich haben sie die Möglichkeit, eigene Kategorien zu erstellen, zu 
bearbeiten und zu löschen, um ihre Ausgaben individuell zu strukturieren und 
zu organisieren.

\subsection{Budgetverwaltung}

Benutzer können Budgets festlegen, um ihre Ausgaben besser zu planen und zu 
kontrollieren. Dabei können sie sowohl ein Gesamtbudget als auch monatliche 
oder kategoriebasierte Budgets definieren. Zusätzlich haben sie die Möglichkeit, 
Budgetlimits jederzeit zu bearbeiten, um flexibel auf Veränderungen zu reagieren.

\subsection{Benachrichtigungen}

Benutzer werden benachrichtigt, wenn sie ihr Budget fast erreichen oder überschreiten. 
Dadurch können sie ihre Ausgaben rechtzeitig anpassen und eine bessere Kontrolle über 
ihre Finanzen behalten.

\subsection{Rechnungsscan}

Benutzer können Rechnungen als Bild oder PDF hochladen und scannen lassen. 
Dabei werden relevante Informationen automatisch erkannt und ausgelesen, 
sodass die Ausgaben einfach und schnell erfasst werden können.

\subsection{Ausgabenprüfung nach Upload}

Nach dem Hochladen einer Rechnung können Benutzer eine Übersicht der 
erkannten Ausgaben einsehen. Diese enthält einzelne Ausgaben, zugewiesene 
Kategorien sowie den Gesamtbetrag. Zusätzlich können Benutzer Details wie Betrag, 
Datum und Kategorie überprüfen und erkannte Kategorien bei Bedarf anpassen, um 
eine korrekte Zuordnung sicherzustellen. Benutzer kann die erfassten Daten bestätigen 
und speichern, nachdem er sie überprüft hat.

\subsection{Erkennung doppelter Rechnungen}

Das System erkennt automatisch mögliche doppelte Rechnungen, um doppelte Buchungen 
zu vermeiden. Benutzer werden entsprechend benachrichtigt, wenn eine Rechnung als 
potenziell doppelt erkannt wird.

\subsection{Manuelle Erfassung von Ausgaben}

Benutzer können Ausgaben manuell eingeben, wenn kein Beleg vorhanden ist. Dabei 
können sie relevante Informationen wie Betrag, Datum und Kategorie angeben, um 
ihre Ausgaben vollständig zu erfassen.

\subsection{Automatische Kategorisierung}

Ausgaben werden automatisch Kategorien zugeordnet, um den Erfassungsprozess zu 
beschleunigen und den Benutzer zu entlasten.

\subsection{Ausgabenanalyse und Visualisierung}

Benutzer können ihre Ausgaben in Form von visuellen Diagrammen analysieren, 
um ihr Ausgabeverhalten besser zu verstehen. Dabei werden wöchentliche und 
monatliche Ausgaben sowie Trends im Zeitverlauf dargestellt. Zusätzlich können 
Benutzer Ausgaben nach Kategorien auswerten und die wichtigsten Kategorien identifizieren, 
um fundierte finanzielle Entscheidungen zu treffen.

\subsection{Personalisierte Empfehlungen}

Benutzer erhalten personalisierte Empfehlungen basierend auf ihrem Ausgabeverhalten, um Einsparpotenziale zu erkennen und ihre Finanzen effizienter zu verwalten.

\end{document}
